\section{Thermodynamic model}\label{sec:theory}
% P:theory-01
The liquid contains nine true species: \COtwo, \MEA, \HtwoO, \MEAH, \MEACOO, \HCOthree, \COthree, \HthreeO and \Hydroxide.
Each equilibrium state is a reactive bubble state of one liquid and one incipient vapor at the measured temperature and apparent composition.
Only the neutral species \COtwo, \MEA and \HtwoO enter the vapor, and the calculated \COtwo partial pressure is \(p_{\COtwo}=y_{\COtwo}P\).
The same ePC-SAFT equation of state describes both phases.

\subsection{Residual Helmholtz energy}
% P:theory-02
The reduced residual Helmholtz energy is the sum of hard-chain, dispersion, association, Debye--H\"uckel and Born contributions \cite{Gross2001,Gross2002,Cameretti2005,Held2014,Figiel2025}:
\begin{equation}
    a^{\mathrm{res}}
    =
    a^{\mathrm{hc}}
    +a^{\mathrm{disp}}
    +a^{\mathrm{assoc}}
    +a^{\mathrm{DH}}
    +a^{\mathrm{Born}} .
    \label{eq:epcsaft-residual-split}
\end{equation}
The hard-chain and dispersion terms follow Gross and Sadowski \cite{Gross2001}, with the dispersion energy of a pair given by
\begin{equation}
    u_{ij}=\sqrt{u_iu_j}\,\left(1-k_{ij}\right).
    \label{eq:dispersion-combining-rule}
\end{equation}
Dispersion between ions of the same charge sign, including an ion with itself, is excluded, while ion--solvent and cation--anion pairs keep \cref{eq:dispersion-combining-rule} \cite{Figiel2025}.
The \MEAH--water interaction depends on temperature,
\begin{equation}
    k_{\MEAH,\HtwoO}(T)=k_{\MEAH,\HtwoO}+b\left(\frac{1}{T}-\frac{1}{T_{\mathrm{ref}}}\right),
    \qquad T_{\mathrm{ref}}=\SI{313.15}{\kelvin},
    \label{eq:kij-temperature}
\end{equation}
and every other \(k_{ij}\) is constant.
Ion hard-chain and Debye--H\"uckel diameters are \(d_i=0.88\,\sigma_i\) \cite[Eq.~4]{Figiel2025}.

% P:theory-03
Water and \MEA carry two association sites (2B), and \COtwo carries two induced-association sites that bond only with water \cite{Kleiner2007,Pabsch2020,Schick2023}.
The \COtwo--water cross-association energy is \SI{1212.85}{\kelvin}, half the water self-association energy, with the water association volume \cite[Table~1, Sect.~2.4.2]{Schick2023}.
\MEA--water cross-association follows the Wolbach--Sandler rule used by Fakouri Baygi and Pahlavanzadeh \cite[Eq.~6]{Baygi2015}.
The association strength carries \(\sigma_{ij}^3\), so their \MEA association volume, reported for a \(d_{ij}^3\) strength \cite[Table~2]{Baygi2015}, was converted from 0.037470 to 0.036287, the constant value that best reproduces the vapor pressure and saturated-liquid density of their model at 293.15--\SI{393.15}{\kelvin}.

\subsection{Born contribution with solvation shell and dielectric saturation}
% P:theory-04
The Born term follows the solvation-shell and dielectric-saturation (SSM+DS) form of Figiel et al.\ \cite[Eqs.~5--10]{Figiel2025}, with its Eq.~7 as corrected in 2026 \cite{Figiel2026Correction}:
\begin{equation}
    a^{\mathrm{Born}}
    =
    -\frac{e^2}{4\pi\varepsilon_0k_{\mathrm{B}}T}
    \sum_{i\in\mathcal{I}}x_iz_i^2
    \left[
        \frac{1-\varepsilon_{\mathrm{bulk}}^{-1}}{D_i}
        +c_{\mathrm{dielectric}}\left(1-\varepsilon_{\mathrm{ion}}^{-1}\right)
        \left(\frac{1}{d_i^{\mathrm{Born}}}-\frac{1}{D_i}\right)
    \right],
    \label{eq:born-ssm-ds}
\end{equation}
\begin{equation}
    D_i=d_i^{\mathrm{Born}}\left[1+c_{\mathrm{shell}}\frac{f_{\mathrm{mix}}-1}{|z_i|}\right],
    \qquad
    f_{\mathrm{mix}}=\frac{\sum_kx_kf_k}{\sum_kx_k}.
    \label{eq:born-shell-diameter}
\end{equation}
Here \(\mathcal{I}\) is the set of ions, \(d_i^{\mathrm{Born}}\) is the Born diameter, \(D_i\) is the outer diameter of the first solvation shell and \(\varepsilon_{\mathrm{ion}}=8\) is the permittivity assigned to the region between the ion surface and that shell.
The selected model uses \(c_{\mathrm{shell}}=c_{\mathrm{dielectric}}=1\), which reproduces the corrected Eq.~7; the original Born comparison uses \(c_{\mathrm{shell}}=c_{\mathrm{dielectric}}=0\), which recovers the original Born expression \cite[Eq.~5]{Figiel2025}.
The water solvation factor is \(f_{\HtwoO}=1.5\) \cite[Table~2]{Figiel2025}, and every other component, including \MEA, \COtwo and all ions, has \(f_k=1\).
Figiel et al.\ restrict the average in \cref{eq:born-shell-diameter} to ions and solvents; here it runs over all nine components, so \(f_{\mathrm{mix}}=1+0.5\,x_{\HtwoO}\).
The shell diameter therefore changes with loading through the water mole fraction, and the Born term contributes to the water activity.

\subsection{Relative permittivity}
% P:theory-05
The bulk permittivity in the Debye--H\"uckel and Born terms is the mass-weighted permittivity of the neutral solvents,
\begin{equation}
    \varepsilon_{\mathrm{bulk}}
    =
    \frac{\sum_{s\in\mathcal{S}}x_sM_s\varepsilon_s}{\sum_{s\in\mathcal{S}}x_sM_s},
    \qquad \mathcal{S}=\{\HtwoO,\MEA\},
    \label{eq:solvent-permittivity}
\end{equation}
which is the salt-free weight-fraction rule of Uyan et al.\ and Figiel et al.\ \cite[Eq.~5]{Uyan2015}\cite[Eq.~12]{Figiel2025}.
The ion-fraction suppression of Figiel et al.\ \cite[Eq.~11]{Figiel2025} is not applied.
The \MEA permittivity is held constant at \(\varepsilon_{\MEA}=32\), an assumed value (\cref{sec:parameters}), and the water permittivity is the software's approximation of the Archer--Wang permittivity,
\begin{equation}
    \varepsilon_{\HtwoO}(T)=78.4104-0.361338\,\theta+7.65556\times10^{-4}\,\theta^{2},
    \qquad \theta=T/\mathrm{K}-298.15 .
    \label{eq:water-permittivity}
\end{equation}

\subsection{Reactions and phase equilibrium}
% P:theory-06
The apparent feed is \SI{1}{\mole} of \MEA, the water amount given by the source composition, and \(\alpha\)~\si{\mole} of \COtwo, where \(\alpha\) is the loading in mol \COtwo per mol total amine.
The water amount follows from the reported \MEA molality where the source gives one and otherwise from the reported \COtwo-free \MEA mass fraction.
The liquid conserves the carbon, hydrogen, nitrogen and oxygen amounts of this feed.
For these nine species the four element balances imply electroneutrality and leave the five independent reactions of \cref{tab:reaction-equilibrium-constants}.
The equilibrium state satisfies equal chemical potentials of \COtwo, \MEA and \HtwoO in the two phases and, for each reaction \(r\),
\begin{equation}
    \sum_i\nu_{ri}\ln a_i=\ln K_r(T),
    \label{eq:activity-reaction-equilibrium}
\end{equation}
where the activities \(a_i\) are evaluated from ePC-SAFT residual chemical potentials on the standard state of each source constant.
Solutes are referenced to infinite dilution in pure water and water to the pure liquid.
A molality-basis constant relates to the corresponding mole-fraction constant by
\begin{equation}
    \ln K_{m,r}=\ln K_{x,r}-\Big(\sum_{i\ne\HtwoO}\nu_{ri}\Big)\ln\!\left(M_{\HtwoO}m^{\circ}\right),
    \label{eq:molality-conversion}
\end{equation}
with \(m^{\circ}=\SI{1}{\mole\per\kilo\gram}\), so \(-\ln(M_{\HtwoO}m^{\circ})=4.0165\).
The R1--R3 source standard states are at system pressure, and the R4 and R5 constants carry their source standard pressure of \SI{100}{\kilo\pascal}.

% P:theory-07
R1--R3 use the correlations of Austgen \cite[Table~V]{Austgen1991},
\begin{equation}
    \ln K_r=A_r+\frac{B_r}{T}+C_r\ln T+D_rT,
    \label{eq:reaction-equilibrium-constant-temperature}
\end{equation}
on the mole-fraction basis, which corresponds to the molality-basis offsets \(\Delta_r\) in \cref{tab:reaction-equilibrium-constants}.
R4 is the carbamate-hydrolysis constant of Tong et al.\ \cite[Table~5]{Tong2012}, converted from Aroua et al.\ \cite{arouaEquilibriumConstantCarbamate1999}, \(\ln K_4=2.151-\SI{1545.3}{\kelvin}/T\) on the aqueous-molality basis.
R5 has the Bates--Pinching form \(\ln K_5=-\ln 10\,(A_5/T+B_5+C_5T)\) on the aqueous-molality basis \cite[Table~5]{Bottinger2008}.
The R2 correlation used here has no separate offset, because its \(A_2\) and \(B_2\) already contain the +4.0165 molality conversion and a fixed constant reaction-enthalpy shift of \SI{-8.2019}{\kilo\joule\per\mole}, applied so that \(\ln K_2(\SI{313.15}{\kelvin})=-14.5050\) is unchanged.
The R5 coefficients used here contain two fixed changes, which \cref{sec:data-methods} describes with the R2 shift.
All five correlations are used over a declared domain of 293.15--\SI{393.15}{\kelvin}; the R4 and R5 source correlations end at \SI{323.15}{\kelvin}, so they are used beyond their source range at higher temperatures.

\begin{table*}[t]
    \centering
    \scriptsize\rmfamily
    \setlength{\tabcolsep}{3pt}
    \renewcommand{\arraystretch}{1.2}
    \caption{Liquid reactions and the recorded equilibrium-constant coefficients of the selected parameter record, with \(T\) in kelvin.
    R1--R3 use \(\ln K_r=A_r+B_r/T+C_r\ln T+D_rT\) on the source mole-fraction basis, and \(\Delta_r\) is the offset of \cref{eq:molality-conversion} to the aqueous-molality basis.
    R2 is recorded directly on the molality basis: its \(A_2\) and \(B_2\) contain the +4.0165 conversion and the inherited \SI{-8.2019}{\kilo\joule\per\mole} enthalpy shift, so it has no separate \(\Delta_2\).
    R4 uses \(\ln K_4=A_4+B_4/T\), and R5 uses \(\ln K_5=-\ln10\,(A_5/T+B_5+C_5T)\), both on the aqueous-molality basis.
    The last column gives the temperature range of the source correlation.}
    \label{tab:reaction-equilibrium-constants}
    \begin{tabularx}{\textwidth}{@{}l>{\raggedright\arraybackslash}p{3.7cm}>{\raggedright\arraybackslash}X>{\raggedright\arraybackslash}p{4.5cm}c@{}}
        \toprule
        ID & Reaction & Source and treatment & Recorded coefficients & Source range (K) \\
        \midrule
        R1 & \(2\HtwoO \rightleftharpoons \HthreeO+\Hydroxide\) & Austgen \cite[Table~V]{Austgen1991}; unchanged & \(A=132.899\), \(B=-13445.9\), \(C=-22.4773\), \(D=0\); \(\Delta=8.0331\) & 273.15--498.15 \\
        R2 & \(\COtwo+2\HtwoO \rightleftharpoons \HCOthree+\HthreeO\) & Austgen \cite[Table~V]{Austgen1991}, \(A=231.465\), \(B=-12092.1\); molality conversion and inherited enthalpy shift included & \(A=232.3314\), \(B=-11105.6400\), \(C=-36.7816\), \(D=0\) & 273.15--498.15 \\
        R3 & \(\HCOthree+\HtwoO \rightleftharpoons \COthree+\HthreeO\) & Austgen \cite[Table~V]{Austgen1991}; unchanged & \(A=216.049\), \(B=-12431.7\), \(C=-35.4819\), \(D=0\); \(\Delta=4.0165\) & 273.15--498.15 \\
        R4 & \(\MEACOO+\HtwoO \rightleftharpoons \MEA+\HCOthree\) & Tong et al.\ \cite[Table~5]{Tong2012}, from Aroua et al.\ \cite{arouaEquilibriumConstantCarbamate1999}; unchanged & \(A=2.151\), \(B=-1545.3\) & 293.15--323.15 \\
        R5 & \(\MEAH+\HtwoO \rightleftharpoons \MEA+\HthreeO\) & Bates--Pinching form \cite[Table~5]{Bottinger2008}, \(A=2677.91\), \(B=0.3869\), \(C=4.277\times10^{-4}\); two inherited changes & \(A=3037.6400\), \(B=-1.0173\), \(C=4.277\times10^{-4}\) & 273.15--323.15 \\
        \bottomrule
    \end{tabularx}
\end{table*}

